\documentclass[12pt,a4paper]{report}

% Поля
\usepackage[
    left=30mm,
    right=15mm,
    top=20mm,
    bottom=20mm
]{geometry}

% Кодировки и русский язык
\usepackage[T2A]{fontenc}
\usepackage[utf8]{inputenc}
\usepackage[russian]{babel}

% Шрифт Tempora (клон Times New Roman с кириллицей)
\usepackage{tempora}

% Математика
\usepackage{amsmath,amssymb}

% Графика
\usepackage{graphicx}

% Таблицы
\usepackage{tabularx}
\usepackage{array}

% Списки (настройка отступов вручную)

% Межстрочный интервал (полуторный)
\linespread{1.5}

% Абзацный отступ
\usepackage{indentfirst}
\setlength{\parindent}{12.5mm}

% Гиперссылки
\usepackage{hyperref}
\hypersetup{colorlinks=true, linkcolor=black, citecolor=black, urlcolor=blue}

% Подсветка кода (listings)
\usepackage{xcolor}
\usepackage{listings}
\lstset{
    language=Python,
    basicstyle=\ttfamily\small,
    backgroundcolor=\color{gray!7},
    frame=single,
    rulecolor=\color{gray!30},
    numbers=left,
    numberstyle=\tiny\color{gray!50},
    numbersep=5pt,
    breaklines=true,
    breakatwhitespace=true,
    keywordstyle=\color{blue!60!black}\bfseries,
    commentstyle=\color{green!40!black},
    stringstyle=\color{red!50!orange},
    showstringspaces=false,
    tabsize=4,
    captionpos=b,
    xleftmargin=20pt,
    framexleftmargin=15pt,
}

% Нумерация страниц
\usepackage{fancyhdr}
\pagestyle{fancy}
\fancyhf{}
\fancyfoot[C]{\thepage}
\renewcommand{\headrulewidth}{0pt}

% Настройка заголовков разделов (вручную)
\makeatletter
\renewcommand{\chapter}{%
  \clearpage\global\@topnum\z@
  \@afterindenttrue
  \secdef\@chapter\@schapter}
\def\@chapter[#1]#2{%
  \ifnum\c@secnumdepth>\m@ne
    \refstepcounter{chapter}%
    \addcontentsline{toc}{chapter}{\protect\numberline{\thechapter}#1}%
  \else
    \addcontentsline{toc}{chapter}{#1}%
  \fi
  \chaptermark{#1}%
  \addtocontents{lof}{\protect\addvspace{10\p@}}%
  \addtocontents{lot}{\protect\addvspace{10\p@}}%
  \@makechapterhead{#2}\@afterheading}
\def\@makechapterhead#1{%
  \vspace*{20\p@}%
  {\parindent\z@\centering\normalfont\bfseries\Large
    Глава \thechapter\par\nobreak
    \vskip 10\p@
    #1\par\nobreak
    \vspace{20\p@}}}
\def\@makeschapterhead#1{%
  \vspace*{20\p@}%
  {\parindent\z@\centering\normalfont\bfseries\Large
    #1\par\nobreak
    \vspace{20\p@}}}
\renewcommand{\section}{%
  \@startsection{section}{1}{\z@}%
  {-15pt}{10pt}%
  {\centering\normalfont\bfseries\large}}
\renewcommand{\subsection}{%
  \@startsection{subsection}{2}{\z@}%
  {-10pt}{5pt}%
  {\normalfont\bfseries\large}}
\makeatother

\begin{document}

\begin{titlepage}
\begin{center}

\textbf{МИНИСТЕРСТВО ЦИФРОВОГО РАЗВИТИЯ, СВЯЗИ И МАССОВЫХ КОММУНИКАЦИЙ РОССИЙСКОЙ ФЕДЕРАЦИИ}

\vspace{0.3em}

\parbox{0.9\textwidth}{\centering\small
Ордена Трудового Красного Знамени федеральное государственное \\
бюджетное образовательное учреждение \\
высшего образования
}

\vspace{0.3em}

\textbf{«Московский технический университет связи и информатики»} \\
(МТУСИ)

\vfill

\textbf{\Large ОТЧЁТ} \\
по производственной практике (проектно-технологической)

\vfill

\begin{flushright}
\begin{minipage}{0.55\textwidth}
\raggedleft
\setlength{\parskip}{0.3em}

\textbf{Выполнил:} \\
Студент группы БПИ2304 \\
Жуков Никита Александрович \\[0.5em]
\makebox[\linewidth]{\hrulefill} \\[0.3em]
(дата, подпись)

\vspace{1em}

\textbf{Проверил руководитель практики:} \\
Жариков Александр Романович \\
Руководитель отдела аналитики и архитектуры \\[0.5em]
\makebox[\linewidth]{\hrulefill} \\[0.3em]
(дата, подпись)
\end{minipage}
\end{flushright}

\vfill

\begin{center}
Москва, 2026
\end{center}

\end{center}
\end{titlepage}

\tableofcontents
\newpage

\chapter*{Введение}
\addcontentsline{toc}{chapter}{Введение}

В современных условиях организации взаимодействия с клиентами и обучающимися сталкиваются с проблемой фрагментации каналов коммуникации. Электронная почта, мессенджеры (Telegram, WhatsApp), образовательные платформы и социальные сети генерируют потоки обращений, которые сложно обрабатывать централизованно. Мониторинг всех этих каналов по отдельности занимает значительное время, и существует высокий риск упустить важное сообщение. В связи с этим разработка «единого окна» коммуникации, объединяющего все каналы в общем интерфейсе оператора, является актуальной задачей в области программной инженерии.

Производственная (проектно-технологическая) практика была направлена на закрепление теоретических знаний, полученных в процессе обучения, и приобретение практических навыков проектирования, разработки и развёртывания современных веб-приложений. Практика проходила в \textbf{ООО «ВИЖНЛАБС»}, где передо мной стояла задача создания омниканального решения для центра обучения компании.

В рамках тестового задания требовалось исследовать существующие на рынке омниканальные платформы (Front, respond.io, Zendesk), сформулировать требования к «единому окну» коммуникации для центра обучения, а затем разработать собственное программное решение, позволяющее принимать и обрабатывать обращения из различных каналов в едином интерфейсе куратора. За основу было принято решение разрабатывать самописное решение, а не использовать готовую SaaS-платформу, поскольку это позволяет точно подстроить систему под реальные процессы центра обучения, особенности используемой образовательной платформы, внутренние роли кураторов и правила обработки обращений, а также сохранить полный контроль над данными и логикой распределения.

\textbf{Целью практики} являлось участие в проектировании и разработке омниканальной платформы для центра обучения, обеспечивающей приём и обработку обращений из различных каналов связи в едином интерфейсе куратора.

В ходе выполнения практики были поставлены и решены следующие \textbf{задачи}:
\begin{enumerate}
    \item Изучить теоретические основы омниканальных решений, проанализировать отличия омниканального подхода от мультиканального.
    \item Спроектировать архитектуру системы, обеспечивающую модульность, расширяемость и слабую связанность компонентов.
    \item Реализовать серверную часть приложения с использованием современных технологий (Python, Litestar, PostgreSQL).
    \item Разработать клиентский интерфейс оператора на базе веб-технологий (React, TypeScript, WebSocket).
    \item Интегрировать каналы связи (Telegram, Email) через единый интерфейс движка сообщений.
    \item Развернуть систему с использованием контейнеризации Docker и провести тестирование ключевых сценариев работы.
\end{enumerate}

Выполнение данных задач позволило получить практические навыки проектирования архитектуры программных систем, разработки серверных приложений на языке Python с использованием асинхронного веб-фреймворка Litestar, работы с реляционными базами данных PostgreSQL и объектным хранилищем MinIO, интеграции внешних API (Telegram Bot API, IMAP/SMTP), а также контейнеризации и развёртывания приложений с использованием Docker.

Настоящий отчёт содержит описание выполненных в ходе практики работ, принятых архитектурных решений, процесса реализации основных компонентов системы и полученных результатов. Отчёт состоит из введения, трёх глав, заключения и списка литературы. В главе 1 рассмотрены теоретические основы омниканальных решений. В главе 2 описана архитектура разработанного программного комплекса. Глава 3 содержит детали реализации ключевых компонентов системы.

\newpage

\chapter{Теоретические основы омниканальных решений}

\section{Понятие омниканальности}

В современном мире цифровых технологий и повсеместного распространения интернета взаимодействие компаний с клиентами претерпело кардинальные изменения. Традиционные модели коммуникации, основанные на использовании одного или нескольких изолированных каналов, уступают место более сложным и интегрированным подходам. Одним из таких подходов является омниканальная стратегия, которая предполагает бесшовную интеграцию всех каналов взаимодействия с клиентом в единую экосистему.

Омниканальность (от лат. \textit{omni}~---~«все» и \textit{channel}~---~«канал») представляет собой стратегический подход к организации взаимодействия с клиентами, при котором все каналы коммуникации (веб-сайт, мобильное приложение, социальные сети, электронная почта, телефон, физические точки продаж и др.) объединены в единую систему. Ключевой особенностью омниканального подхода является то, что клиент воспринимается как единое целое независимо от того, через какой канал он взаимодействует с компанией, а все каналы обмениваются данными в режиме реального времени.

Важно различать понятия мультиканальности и омниканальности. Мультиканальный подход предполагает использование нескольких каналов коммуникации, однако эти каналы часто существуют изолированно друг от друга. Например, клиент может оформить заказ на веб-сайте, но при обращении в контакт-центр ему придётся заново сообщать всю информацию. В омниканальной модели такая ситуация исключена: история взаимодействия клиента доступна во всех каналах одновременно, что обеспечивает непрерывность клиентского опыта.

\section{Эволюция каналов взаимодействия с клиентами}

Развитие каналов взаимодействия с клиентами прошло несколько этапов. На начальном этапе компании использовали исключительно офлайн-каналы: личные встречи, телефонные звонки, почтовые рассылки. С развитием интернета появились онлайн-каналы: веб-сайты, электронная почта, чаты. Вслед за этим возникла необходимость интеграции различных каналов, что привело к появлению мультиканального подхода.

\begin{table}[ht]
\centering
\caption{Эволюция подходов к взаимодействию с клиентами}
\label{tab:evolution}
\begin{tabularx}{\textwidth}{|X|X|X|}
\hline
\textbf{Характеристика} & \textbf{Мультиканальный подход} & \textbf{Омниканальный подход} \\ \hline
Интеграция каналов & Частичная или отсутствует & Полная интеграция \\ \hline
Единая история клиента & Отсутствует & Едина для всех каналов \\ \hline
Качество обслуживания & Различается по каналам & Единое на всех каналах \\ \hline
Аналитика & Разрозненная & Консолидированная \\ \hline
\end{tabularx}
\end{table}

Современный этап развития характеризуется переходом к омниканальности, что обусловлено несколькими факторами. Во-первых, изменилось поведение потребителей: современный клиент использует в среднем 4--6 различных каналов при взаимодействии с компанией. Во-вторых, развитие технологий позволило обеспечить техническую возможность интеграции разнородных систем. В-третьих, конкурентная среда вынуждает компании искать новые способы привлечения и удержания клиентов.

\section{Ключевые компоненты омниканальной архитектуры}

Омниканальная архитектура включает в себя несколько взаимосвязанных компонентов:

\begin{enumerate}
    \item \textbf{Единая платформа управления взаимодействием}~---~центральный элемент архитектуры, обеспечивающий сбор, хранение и обработку данных о взаимодействиях с клиентами из всех каналов.

    \item \textbf{Интеграционный слой}~---~совокупность программных интерфейсов (API) и middleware-решений, обеспечивающих обмен данными между различными системами.

    \item \textbf{Система управления контентом}~---~обеспечивает единообразное представление информации на всех каналах.

    \item \textbf{Аналитическая платформа}~---~собирает и анализирует данные о поведении клиентов, позволяя принимать обоснованные решения.

    \item \textbf{Система управления взаимоотношениями с клиентами (CRM)}~---~хранит полную историю взаимодействия с каждым клиентом.
\end{enumerate}

Особое значение в омниканальной архитектуре имеет качество данных. Поскольку информация поступает из множества источников, критически важным является обеспечение её непротиворечивости и актуальности.

\section{Преимущества внедрения омниканальных решений}

Внедрение омниканального подхода предоставляет компаниям ряд существенных преимуществ:

\begin{itemize}
    \item \textbf{Повышение лояльности клиентов.} Бесшовный клиентский опыт способствует росту удовлетворённости и, как следствие, увеличению показателя удержания клиентов (retention rate).

    \item \textbf{Увеличение средней выручки на клиента.} Интеграция каналов открывает возможности для кросс-канальных продаж и персонализированных предложений.

    \item \textbf{Оптимизация операционных затрат.} Автоматизация процессов и единая аналитика позволяют сократить издержки на обслуживание клиентов.

    \item \textbf{Повышение эффективности маркетинга.} Консолидированные данные о клиентах обеспечивают более точное таргетирование и персонализацию коммуникаций.
\end{itemize}

Исследования показывают, что компании, успешно внедрившие омниканальную стратегию, демонстрируют на 10--15\% более высокие показатели удовлетворённости клиентов и на 20--30\% больший рост выручки по сравнению с компаниями, использующими мультиканальный подход.

\section{Технологические аспекты реализации}

Реализация омниканального подхода требует использования современных технологических решений. Ключевую роль играют облачные технологии, обеспечивающие масштабируемость и доступность систем. Микросервисная архитектура позволяет гибко наращивать функциональность и интегрировать новые каналы.

В области управления данными широкое применение находят технологии Big Data и системы класса Customer Data Platform (CDP), которые обеспечивают сбор и унификацию данных из различных источников. Машинное обучение используется для персонализации взаимодействий и прогнозирования поведения клиентов.

Безопасность данных является критически важным аспектом, особенно в контексте требований законодательства о защите персональных данных. Омниканальные решения должны обеспечивать соответствие требованиям Федерального закона №~152-ФЗ «О персональных данных» и других нормативных актов.

\section{Отраслевые особенности применения}

Омниканальные решения находят применение в различных отраслях экономики. В розничной торговле они позволяют объединить онлайн- и офлайн-каналы продаж, обеспечивая такие сценарии, как «клик-энд-коллект» (заказ онлайн с самовывозом в магазине) и возврат товаров, приобретённых онлайн, в физических точках продаж.

В банковском секторе омниканальность проявляется в интеграции мобильного банкинга, интернет-банка, отделений и контакт-центра. Клиент может начать оформление продукта в мобильном приложении и завершить его в отделении без повторного ввода данных.

В сфере телекоммуникаций омниканальный подход обеспечивает единое обслуживание клиентов через личный кабинет, мобильное приложение, чаты поддержки и центры продаж.

\section*{Выводы по главе 1}

Омниканальные решения представляют собой современный подход к организации взаимодействия с клиентами, обеспечивающий бесшовную интеграцию всех каналов коммуникации. Внедрение такого подхода требует значительных инвестиций в технологии и организационные изменения, однако обеспечивает существенные конкурентные преимущества за счёт повышения качества клиентского опыта и эффективности бизнес-процессов.

\newpage
\chapter{Архитектура программного комплекса омниканального решения}

\section{Общая архитектура системы}

Разработанный программный комплекс представляет собой омниканальную платформу для взаимодействия с клиентами, объединяющую несколько каналов связи в едином интерфейсе. Система построена по принципу \textbf{Clean Architecture} (чистая архитектура) с монолитным развёртыванием и состоит из трёх основных уровней: клиентское приложение (фронтенд), серверное приложение (бэкенд) и внешние сервисы.

Общая структурная схема системы представлена ниже:

\medskip
\begin{center}
\footnotesize
\begin{tabular}{c}
\fbox{\parbox{0.3\textwidth}{\centering\vspace{2pt}Клиентское приложение\\(\textbf{Frontend})\vspace{2pt}}} \\
$\uparrow$ HTTP/WS \\
\fbox{\parbox{0.5\textwidth}{\centering\vspace{2pt}\textbf{Серверное приложение (Backend)}\\\vspace{2pt}}} \\
$\downarrow$ $\downarrow$ $\downarrow$ \\
\fbox{\parbox{0.2\textwidth}{\centering\vspace{2pt}\small PostgreSQL\vspace{2pt}}}
\hspace{1em}
\fbox{\parbox{0.2\textwidth}{\centering\vspace{2pt}\small MinIO (S3)\vspace{2pt}}}
\hspace{1em}
\fbox{\parbox{0.22\textwidth}{\centering\vspace{2pt}\small Telegram Bot API\vspace{2pt}}}
\hspace{1em}
\fbox{\parbox{0.18\textwidth}{\centering\vspace{2pt}\small IMAP/SMTP\vspace{2pt}}}
\end{tabular}
\end{center}
\medskip

Основные компоненты системы развёртываются в среде контейнеризации Docker, что обеспечивает изолированность окружений и простоту масштабирования. Конфигурация взаимодействия компонентов задаётся через переменные окружения.

\section{Серверная часть (Backend)}

\subsection{Архитектурный подход}

Серверная часть реализована на языке \textbf{Python 3.14+} с использованием асинхронного веб-фреймворка \textbf{Litestar 2.x}. В основе архитектуры лежит принцип разделения ответственности (Clean Architecture), который предполагает чёткое разделение кода на три слоя:

\begin{itemize}
    \item \textbf{Core (ядро)}~---~содержит доменные модели, бизнес-логику и порты (интерфейсы) для взаимодействия с внешними системами.
    \item \textbf{API}~---~предоставляет REST-эндпоинты и WebSocket для взаимодействия с клиентским приложением.
    \item \textbf{Adapters (адаптеры)}~---~реализуют порты ядра для конкретных технологий (PostgreSQL, Telegram, Email, MinIO).
\end{itemize}

Такая архитектура обеспечивает слабую связанность компонентов и возможность замены любой реализации без изменения бизнес-логики.

\subsection{Схема слоёв бэкенда}

\begin{center}
\footnotesize
\begin{tabular}{c}
\fbox{\parbox{0.55\textwidth}{\centering\vspace{2pt}\textbf{API Layer} --- REST-контроллеры, схемы DTO, обработка исключений\vspace{2pt}}} \\
$\uparrow$ \\
\fbox{\parbox{0.55\textwidth}{\centering\vspace{2pt}\textbf{Core Layer} --- Доменные модели, сервисы (бизнес-логика), порты\vspace{2pt}}} \\
$\uparrow$ \\
\fbox{\parbox{0.55\textwidth}{\centering\vspace{2pt}\textbf{Adapters Layer} --- Реализации: БД, Telegram, Email, S3, файловое хранилище\vspace{2pt}}}
\end{tabular}
\end{center}

\subsection{Технологический стек бэкенда}

\begin{itemize}
    \item \textbf{Веб-фреймворк}: Litestar 2.x (асинхронный, поддержка OpenAPI, WebSocket)
    \item \textbf{ORM}: SQLAlchemy 2.x (асинхронный режим)
    \item \textbf{База данных}: PostgreSQL 17 (через asyncpg)
    \item \textbf{DI-контейнер}: punq (внедрение зависимостей)
    \item \textbf{Аутентификация}: JWT (HS256) + bcrypt
    \item \textbf{Файловое хранилище}: MinIO (S3-совместимое объектное хранилище)
    \item \textbf{Telegram}: python-telegram-bot 22.x (long-polling)
    \item \textbf{Email}: IMAP (входящие), SMTP (исходящие)
\end{itemize}

\subsection{API-эндпоинты}

Серверная часть предоставляет следующие группы эндпоинтов:

\begin{itemize}
    \item \textbf{Сообщения}: получение списка, просмотр, отправка ответа, загрузка файлов
    \item \textbf{Клиенты}: создание, просмотр, редактирование, поиск по каналу
    \item \textbf{Кураторы}: управление учётными записями операторов поддержки
    \item \textbf{Назначения}: привязка куратора к сообщению, передача между кураторами, история
    \item \textbf{Аутентификация}: регистрация, вход, получение текущего пользователя
    \item \textbf{WebSocket}: уведомления о новых сообщениях и переназначениях в реальном времени
\end{itemize}

Все REST-ответы возвращаются в формате JSON. При возникновении ошибок используется формат Problem Details (RFC~9457).

\section{Клиентская часть (Frontend)}

Клиентское приложение представляет собой одностраничное приложение (SPA), реализованное с использованием современного стека веб-технологий:

\begin{itemize}
    \item \textbf{Фреймворк}: React 19
    \item \textbf{Язык}: TypeScript 5.8 (строгая типизация)
    \item \textbf{Сборщик}: Vite 6
    \item \textbf{Стилизация}: Tailwind CSS 3 (адаптивный дизайн, поддержка тёмной темы)
\end{itemize}

\subsection{Структура пользовательского интерфейса}

Интерфейс состоит из двух основных панелей:

\begin{itemize}
    \item \textbf{Левая панель} (320\,px): отображает список контактов с возможностью поиска, сортировки и фильтрации по каналу связи. Для каждого контакта отображается аватар, имя, последнее сообщение и количество непрочитанных сообщений.
    \item \textbf{Правая панель}: содержит область просмотра сообщений (история диалога) и форму ввода ответа. Форма ввода адаптируется под тип канала: для чатов~---~поле ввода с прикреплением файлов, для email~---~отдельная форма с полями «Кому», «Тема» и «Тело письма».
\end{itemize}

\subsection{Взаимодействие с сервером}

Обмен данными между клиентом и сервером осуществляется по протоколу HTTP (REST API) для основных операций и через WebSocket для получения уведомлений в реальном времени. При загрузке страницы устанавливается WebSocket-соединение с автоматическим переподключением при разрыве (с задержкой 3~секунды).

Аутентификация выполняется по JWT-токену, который сохраняется в \texttt{localStorage} браузера и передаётся в заголовке \texttt{Authorization: Bearer <token>} при каждом запросе.

\section{Схема базы данных}

Система использует реляционную базу данных PostgreSQL 17, содержащую следующие основные таблицы:

\begin{itemize}
    \item \textbf{messages}~---~хранит все сообщения из всех каналов. Ключевые поля: идентификатор (формат \texttt{канал:оригинальный\_id}), тип канала, идентификатор отправителя, содержимое, тип сообщения (текст, фото, документ и т.д.), идентификатор родительского сообщения (для тредов), назначенный куратор, временная метка.
    \item \textbf{senders}~---~отправители сообщений (физические или юридические лица). Содержит имя и связанные идентичности.
    \item \textbf{sender\_identities}~---~привязка отправителя к конкретному каналу (например, Telegram ID или email-адрес).
    \item \textbf{clients}~---~клиенты, с которыми ведётся диалог. Содержит имя, телефон, email, ссылку на аватар, метаданные.
    \item \textbf{client\_channels}~---~привязка клиента к каналам связи (канал + внешний идентификатор + имя пользователя).
    \item \textbf{curators}~---~операторы поддержки (кураторы). Содержит ФИО, логин, email, роль, статус, хеш пароля.
    \item \textbf{assignment\_history}~---~история назначения кураторов на сообщения (кто, от кого, кому, причина, временная метка).
\end{itemize}

\begin{center}
\begin{table}[ht]
\centering
\caption{Сводная таблица базы данных}
\label{tab:db_schema}
\begin{tabularx}{\textwidth}{|X|c|X|}
\hline
\textbf{Таблица} & \textbf{Кол-во полей} & \textbf{Назначение} \\ \hline
messages & 11 & Все сообщения из всех каналов \\ \hline
senders & 3 & Отправители сообщений \\ \hline
sender\_identities & 4 & Привязка отправителей к каналам \\ \hline
clients & 7 & Клиенты системы \\ \hline
client\_channels & 5 & Каналы клиентов \\ \hline
curators & 9 & Операторы поддержки \\ \hline
assignment\_history & 7 & История назначений \\ \hline
\end{tabularx}
\end{table}
\end{center}

\section{Модуль интеграции каналов связи}

Ключевым компонентом системы является модуль интеграции каналов, реализованный с использованием паттерна \textbf{Абстрактная фабрика} (Abstract Factory). Каждый канал представлен отдельным классом-движком (Engine), реализующим единый интерфейс для получения и отправки сообщений.

\begin{center}
\footnotesize
\begin{tabular}{c}
\fbox{\parbox{0.55\textwidth}{\centering\vspace{2pt}\textbf{BaseEngine} (интерфейс)\\\small get\_incoming(), send\_message()\vspace{2pt}}} \\
$\swarrow$ \hspace{2em} $\downarrow$ \hspace{2em} $\searrow$ \\
\fbox{\parbox{0.17\textwidth}{\centering\vspace{2pt}\small\textbf{Telegram}\\Engine\vspace{2pt}}}
\hspace{1em}
\fbox{\parbox{0.17\textwidth}{\centering\vspace{2pt}\small\textbf{Email}\\Engine\vspace{2pt}}}
\hspace{1em}
\fbox{\parbox{0.17\textwidth}{\centering\vspace{2pt}\small\textbf{SMS/VK}\\Engine (план)\vspace{2pt}}}
\end{tabular}
\end{center}

\subsection{Модуль Telegram}

Интеграция с мессенджером Telegram реализована через официальное API с использованием библиотеки \texttt{python-telegram-bot}. Модуль работает в режиме long-polling, периодически опрашивая сервер Telegram на наличие новых сообщений. Поддерживаются следующие типы контента:

\begin{itemize}
    \item Текстовые сообщения
    \item Фотографии (загружаются в MinIO, фронтенду передаётся presigned-ссылка)
    \item Документы (аналогичная обработка)
    \item Стикеры (поддерживаются форматы .webp, .webm, .tgs)
    \item Анимации (GIF)
    \item Видео, аудио и голосовые сообщения
\end{itemize}

При поступлении нового сообщения модуль автоматически создаёт или обновляет запись клиента в базе данных, сохраняет сообщение и уведомляет подключённых операторов через WebSocket.

\subsection{Модуль Email}

Интеграция с электронной почтой реализована по протоколам IMAP (для получения входящих сообщений) и SMTP (для отправки ответов). Поддерживаются как текстовые, так и HTML-письма. На фронтенде HTML-часть письма отображается в iframe-контейнере.

Цикл обработки email-сообщений:
\begin{enumerate}
    \item Периодический опрос IMAP-сервера на предмет новых непрочитанных писем.
    \item Извлечение текстового и HTML-содержимого, декодирование темы и заголовков.
    \item Сохранение в базе данных в формате, едином для всех каналов.
    \item Отправка ответа через SMTP с сохранением темы письма (префикс <<Re:>>).
\end{enumerate}

\subsection{Планируемые каналы}

Архитектура системы предусматривает возможность подключения дополнительных каналов без изменения существующего кода. В настоящий момент в кодовой базе предусмотрена поддержка (включена в перечисление типов каналов, но не реализована):

\begin{itemize}
    \item \textbf{SMS}~---~планируется через интеграцию с SMS-шлюзом
    \item \textbf{VK}~---~планируется через VK API
\end{itemize}

Подключение нового канала сводится к реализации класса-движка, наследующего интерфейс BaseEngine, и регистрации его фабрики в общем реестре.

\section{Взаимодействие компонентов}

\subsection{Получение входящего сообщения}

Процесс получения и обработки входящего сообщения включает следующие шаги:

\begin{enumerate}
    \item \textbf{Пинг канала}: компонент PollingOrchestrator запускает асинхронную задачу для каждого активного движка канала.
    \item \textbf{Получение}: движок канала (Telegram, Email) получает новое сообщение от внешнего API.
    \item \textbf{Создание/обновление клиента}: система проверяет, существует ли клиент с указанным внешним идентификатором. Если нет~---~создаётся новый клиент.
    \item \textbf{Сохранение}: сообщение сохраняется в таблицу \texttt{messages} через репозиторий.
    \item \textbf{Уведомление}: через WebSocket всем подключённым операторам отправляется событие \texttt{new\_message}.
    \item \textbf{Отображение}: фронтенд получает событие и добавляет сообщение в интерфейс соответствующего оператора.
\end{enumerate}

\subsection{Отправка ответа}

\begin{enumerate}
    \item Оператор вводит текст ответа в интерфейсе и нажимает <<Отправить>>.
    \item Фронтенд отправляет POST-запрос на \texttt{/api/messages/reply}.
    \item Бэкенд создаёт объект сообщения-ответа со ссылкой на родительское сообщение (\texttt{parent\_id}).
    \item Движок соответствующего канала отправляет сообщение через внешнее API (Telegram~---~в чат, Email~---~на почтовый адрес).
    \item Ответ сохраняется в базу данных и транслируется через WebSocket.
\end{enumerate}

\subsection{Назначение куратора}

\begin{enumerate}
    \item Оператор выбирает сообщение и назначает ответственного куратора (себя или другого оператора).
    \item Фронтенд отправляет запрос на \texttt{POST /api/messages/\{id\}/assign}.
    \item Бэкенд обновляет поле \texttt{curator\_id} в записи сообщения и создаёт запись в \texttt{assignment\_history}.
    \item Система уведомляет всех операторов через WebSocket о назначении.
\end{enumerate}

\newpage
\section*{Выводы по главе 2}

Разработанный программный комплекс реализует архитектуру, обеспечивающую надёжную и масштабируемую работу омниканальной платформы. Использование Clean Architecture позволяет изолировать бизнес-логику от технических деталей реализации, что упрощает поддержку и развитие системы. Модульная архитектура интеграции каналов обеспечивает возможность подключения новых каналов связи с минимальными трудозатратами. Применение современных технологий (Python 3.14, Litestar, React 19, PostgreSQL 17, WebSocket) гарантирует высокую производительность и отзывчивость интерфейса.

\newpage
\chapter{Реализация программного комплекса}

В данной главе рассматриваются ключевые аспекты реализации системы, включая доменную модель, бизнес-логику, адаптеры каналов связи, репозитории для работы с базой данных, а также конфигурацию и развёртывание.

\section{Доменная модель}

Центральным элементом предметной области является класс \texttt{Message}, представляющий сообщение из любого канала связи. Модель реализована в виде неизменяемого dataclass-а с использованием \texttt{slots=True} для оптимизации потребления памяти:

\begin{lstlisting}[language=Python]
@dataclass(slots=True)
class Message:
    id: str
    channel: str
    sender_id: str
    content: str
    timestamp: datetime
    metadata: Dict[str, Any]
    message_type: MessageType = MessageType.TEXT
    recipient: Optional[str] = None
    subject: Optional[str] = None
    parent_id: Optional[str] = None
    curator_id: Optional[str] = None

    def to_dict(self) -> Dict[str, Any]:
        return {
            "id": self.id,
            "channel": self.channel,
            "sender_id": self.sender_id,
            "content": self.content,
            "timestamp": self.timestamp.isoformat(),
            "message_type": self.message_type.value,
            "metadata": self.metadata,
            "recipient": self.recipient,
            "subject": self.subject,
            "parent_id": self.parent_id,
            "curator_id": self.curator_id,
        }
\end{lstlisting}

Поле \texttt{channel} принимает одно из значений перечисления \texttt{ChannelType}, которое определяет поддерживаемые каналы связи:

\begin{lstlisting}[language=Python]
class ChannelType(StrEnum):
    TELEGRAM = "telegram"
    EMAIL = "email"
    SMS = "sms"
    VK = "vk"
\end{lstlisting}

Тип вложенного контента (текст, фото, документ, стикер и т.д.) определяется перечислением \texttt{MessageType}. Такой подход обеспечивает единообразное представление разнородных сообщений из всех каналов, что является ключевым требованием омниканальной системы.

\section{Порты (интерфейсы)}

Для обеспечения слабой связанности компонентов каждый модуль зависит от абстракции (порта), а не от конкретной реализации. Интерфейс движка канала связи определён следующим образом:

\begin{lstlisting}[language=Python]
class MessageEngine(ABC):
    @abstractmethod
    async def start(self) -> None:
        ...

    @abstractmethod
    async def send_message(
        self, recipient: str, content: str, **kwargs
    ) -> str:
        ...

    @abstractmethod
    def get_incoming(self) -> AsyncIterator[Message]:
        ...

    @property
    @abstractmethod
    def channel_type(self) -> str:
        ...
\end{lstlisting}

Метод \texttt{get\_incoming()} возвращает асинхронный итератор, что позволяет эффективно обрабатывать поток входящих сообщений без блокировки основного потока выполнения. Метод \texttt{send\_message()} возвращает идентификатор отправленного сообщения для последующей синхронизации состояния.

Аналогичным образом определены порты для репозиториев сообщений, клиентов, кураторов и шлюза базы данных.

\section{Бизнес-логика}

\subsection{Сервис сообщений}

\texttt{MessageService} реализует сценарии использования, связанные с управлением сообщениями. Основные операции делегируются репозиторию через интерфейс \texttt{MessageRepository}:

\begin{lstlisting}[language=Python]
class MessageService:
    def __init__(self, repository: MessageRepository):
        self._repository = repository

    async def save_message(self, message: Message) -> None:
        await self._repository.save_message(message)

    async def get_messages(
        self,
        channel: Optional[str] = None,
        sender_id: Optional[str] = None,
        curator_id: Optional[str] = None,
        limit: int = 100,
        offset: int = 0,
    ) -> List[Message]:
        return await self._repository.get_messages(
            channel=channel, sender_id=sender_id,
            curator_id=curator_id, limit=limit, offset=offset,
        )

    async def reply_to_message(
        self, message_id: str, content: str
    ) -> Message:
        original = await self._repository.get_message_by_id(message_id)
        if not original:
            raise MessageNotFoundError(message_id)
        return original
\end{lstlisting}

Метод \texttt{reply\_to\_message()} проверяет существование исходного сообщения и в случае его отсутствия генерирует исключение \texttt{MessageNotFoundError}, которое в дальнейшем преобразуется в HTTP-ответ с кодом 404 в формате Problem Details (RFC~9457).

\subsection{Оркестратор опроса каналов}

\texttt{PollingOrchestrator} отвечает за непрерывный сбор входящих сообщений из всех зарегистрированных каналов. Класс реализует паттерн <<Наблюдатель>> (Observer) и использует асинхронные задачи для параллельного опроса движков:

\begin{lstlisting}[language=Python]
class PollingOrchestrator(AbstractPollingService):
    def __init__(self, message_service, client_service):
        self._message_service = message_service
        self._client_service = client_service
        self._engines: List[MessageEngine] = []
        self._tasks: List[asyncio.Task] = []

    def register_engine(self, engine: MessageEngine) -> None:
        self._engines.append(engine)

    async def start_polling(self) -> None:
        self._running = True
        for engine in self._engines:
            await engine.start()
            task = asyncio.create_task(
                self._poll_engine(engine)
            )
            self._tasks.append(task)

    async def send_reply(self, channel, recipient,
                         content, **kwargs) -> str | None:
        for engine in self._engines:
            if engine.channel_type == channel:
                return await engine.send_message(
                    recipient=recipient, content=content, **kwargs
                )
        return None
\end{lstlisting}

Основной цикл обработки реализован в методе \texttt{\_poll\_engine()}. Для каждого полученного сообщения система автоматически выполняет следующие действия:

\begin{enumerate}
    \item Извлечение информации об отправителе из метаданных канала.
    \item Создание или обновление записи клиента через \texttt{ClientService}.
    \item Сохранение сообщения в базу данных.
    \item Трансляция события \texttt{new\_message} всем подключённым операторам через WebSocket.
\end{enumerate}

Извлечение информации об отправителе учитывает особенности каждого канала: для Telegram извлекаются \texttt{username} и \texttt{first\_name}, для email~---~имя и адрес из заголовка <<From>>.

\section{Адаптеры каналов связи}

\subsection{Telegram}

Адаптер Telegram реализован с использованием библиотеки \texttt{python-telegram-bot}. Класс \texttt{TelegramEngine} обрабатывает различные типы вложений, загружая медиафайлы в объектное хранилище MinIO и формируя presigned-ссылки для доступа со стороны фронтенда.

Метод отправки сообщения поддерживает как текстовые сообщения, так и сообщения с вложениями:

\begin{lstlisting}[language=Python]
async def send_message(
    self, recipient: str, content: str, **kwargs
) -> str:
    if "file_url" in kwargs:
        with aiohttp.ClientSession() as session:
            async with session.get(kwargs["file_url"]) as resp:
                file_data = await resp.read()
        file = InputFile(BytesIO(file_data),
                         filename=kwargs.get("filename", "file"))
        sent = await self._bot.send_document(
            chat_id=recipient, document=file,
            caption=content or None
        )
    else:
        sent = await self._bot.send_message(
            chat_id=recipient, text=content
        )
    return str(sent.id)
\end{lstlisting}

Для загрузки вложений в S3-хранилище реализованы вспомогательные методы:

\begin{lstlisting}[language=Python]
async def _upload_photo_to_s3(
    self, photo: PhotoSize, message_id: str
) -> Optional[str]:
    file = await photo.get_file()
    file_bytes = BytesIO()
    await file.download_to_memory(file_bytes)
    file_bytes.seek(0)
    url = await self._s3.upload_file(
        bucket="attachments",
        key=f"{message_id}/photo.jpg",
        data=file_bytes.getvalue(),
    )
    return url
\end{lstlisting}

\subsection{Email}

Адаптер электронной почты использует встроенные библиотеки Python \texttt{imaplib} (получение) и \texttt{smtplib} (отправка), работая через \texttt{asyncio.run\_in\_executor()} для совместимости с асинхронной архитектурой. Цикл опроса IMAP-сервера:

\begin{lstlisting}[language=Python]
async def get_incoming(self) -> AsyncIterator[Message]:
    while self._running:
        messages = await asyncio.get_event_loop().run_in_executor(
            None, self._fetch_new_messages
        )
        for msg in messages:
            yield msg
        await asyncio.sleep(self._poll_interval)
\end{lstlisting}

Извлечение содержимого письма обрабатывает MIME-структуру, поддерживая как текстовые, так и HTML-части:

\begin{lstlisting}[language=Python]
def _extract_body(self, msg) -> tuple[str, str | None]:
    text_body = ""
    html_body = None
    if msg.is_multipart():
        for part in msg.walk():
            content_type = part.get_content_type()
            if content_type == "text/plain":
                text_body += self._decode_part(part)
            elif content_type == "text/html":
                html_body = self._decode_part(part)
    else:
        text_body = self._decode_part(msg)
    return text_body, html_body
\end{lstlisting}

\subsection{Фабрика движков}

Для создания экземпляров движков используется паттерн <<Абстрактная фабрика>>. Каждый тип канала имеет свою фабрику, которая регистрируется в общем реестре:

\begin{lstlisting}[language=Python]
class TelegramEngineFactory(EngineFactory):
    def create_engine(self, config: dict) -> MessageEngine:
        return TelegramEngine(token=config["token"], s3=self._s3)

    def get_supported_channel(self) -> str:
        return ChannelType.TELEGRAM


class EngineAbstractFactory:
    def __init__(self):
        self._factories: Dict[str, EngineFactory] = {}

    def register_factory(self, factory: EngineFactory) -> None:
        channel = factory.get_supported_channel()
        self._factories[channel] = factory

    def create_engine(self, channel_type, config) -> MessageEngine:
        if channel_type not in self._factories:
            raise ValueError(
                f"No factory for channel: {channel_type}"
            )
        return self._factories[channel_type].create_engine(config)
\end{lstlisting}

Такой подход позволяет добавлять новые каналы связи простой реализацией нового движка и его фабрики, без модификации существующего кода.

\section{Репозиторий и ORM}

Для работы с базой данных используется SQLAlchemy 2.x в асинхронном режиме. Модель сообщения в ORM представляет собой отображение на таблицу \texttt{messages}:

\begin{lstlisting}[language=Python]
class MessageModel(Base):
    __tablename__ = "messages"

    id = Column(String, primary_key=True)
    channel = Column(String(50), nullable=False, index=True)
    sender_id = Column(String, ForeignKey("senders.id",
                                          ondelete="SET NULL"),
                       nullable=True, index=True)
    content = Column(Text, nullable=False)
    message_type = Column(String(20), nullable=False,
                          server_default="text")
    subject = Column(String(500), nullable=True)
    parent_id = Column(String, nullable=True, index=True)
    curator_id = Column(String, nullable=True, index=True)
    timestamp = Column(DateTime(timezone=True), nullable=False)
    metadata_ = Column("metadata", JSON, nullable=True)
    created_at = Column(DateTime(timezone=True),
                        server_default=func.now())
\end{lstlisting}

Репозиторий реализует интерфейс \texttt{MessageRepository}, инкапсулируя детали выполнения SQL-запросов. Методы репозитория работают с доменными объектами, а не с ORM-моделями, что обеспечивает независимость бизнес-логики от выбранной СУБД.

\section{Контейнер внедрения зависимостей}

Для связывания компонентов используется DI-контейнер \texttt{punq}. Все зависимости регистрируются в единственной точке конфигурации:

\begin{lstlisting}[language=Python]
def configure_container() -> punq.Container:
    container = punq.Container()

    container.register(EngineAbstractFactory,
                       instance=EngineAbstractFactory())
    container.register(S3Interface, instance=MinioGateway(...))
    container.register(DatabaseGateway,
                       instance=PostgresDatabaseGateway(...))
    container.register(MessageRepository,
                       instance=PostgresMessageRepository(...))
    container.register(MessageService)
    container.register(ClientService)
    container.register(AbstractPollingService,
                       instance=PollingOrchestrator(...))
    container.register(AuthService)
    container.register(ValidationService)

    return container
\end{lstlisting}

Контейнер обеспечивает автоматическое разрешение зависимостей: если сервис требует \texttt{MessageRepository}, контейнер внедряет соответствующий экземпляр без необходимости ручного конструирования.

Точка входа в приложение~---~функция создания приложения Litestar:

\begin{lstlisting}[language=Python]
@asynccontextmanager
async def lifespan(app: Litestar):
    container = configure_container()
    db = container.resolve(DatabaseGateway)
    await db.init()
    polling = container.resolve(AbstractPollingService)
    await polling.start_polling()
    yield
    await polling.stop_polling()
    await db.close()


app = Litestar(
    route_handlers=[...],
    lifespan=[lifespan],
    websocket_handler=websocket_handler,
    on_app_init=[...],
)
\end{lstlisting}

\section{Конфигурация и развёртывание}

Система развёртывается с использованием Docker Compose, что обеспечивает изолированное окружение для всех компонентов:

\begin{lstlisting}[language={},frame=single,backgroundcolor=\color{gray!7},numbers=none]
services:
  postgres:
    image: postgres:17
    environment:
      POSTGRES_USER: user
      POSTGRES_PASSWORD: password
      POSTGRES_DB: omnichannel
    healthcheck:
      test: ["CMD-SHELL", "pg_isready -U user -d omnichannel"]
    volumes:
      - postgres_data:/var/lib/postgresql/data

  minio:
    image: minio/minio
    command: server /data --console-address ":9001"
    environment:
      MINIO_ROOT_USER: minioadmin
      MINIO_ROOT_PASSWORD: minioadmin
    volumes:
      - minio_data:/data

  app:
    build: .
    env_file: .env
    ports:
      - "8044:8044"
    depends_on:
      postgres: { condition: service_healthy }
\end{lstlisting}

Конфигурация приложения задаётся через переменные окружения (файл \texttt{.env}), что позволяет гибко настраивать параметры подключения к базам данных, токены каналов и настройки безопасности без изменения кода.

\section{Тестирование}

Покрытие кода модульными и интеграционными тестами реализовано с использованием \texttt{pytest} совместно с \texttt{pytest-asyncio} для асинхронных тестов. Пример теста сервиса опроса каналов:

\begin{lstlisting}[language=Python]
async def test_poll_engine_saves_messages(
    mock_engine, mock_message_service, mock_client_service
):
    orchestrator = PollingOrchestrator(
        message_service=mock_message_service,
        client_service=mock_client_service,
    )
    orchestrator.register_engine(mock_engine)

    mock_engine.set_messages([...])
    await orchestrator._poll_engine(mock_engine)

    assert mock_message_service.save_message.called
\end{lstlisting}

Тесты репозиториев используют Docker-контейнер с PostgreSQL (через \texttt{pytest-postgresql} или тестовую БД), что гарантирует корректность SQL-запросов и целостность данных. На момент написания отчёта тестовое покрытие включает тестирование всех сервисов, адаптеров, API-маршрутов и обработчиков исключений.

\newpage
\section*{Выводы по главе 3}

В ходе реализации программного комплекса были применены современные подходы к разработке: Clean Architecture для разделения ответственности, паттерны Abstract Factory и Observer для гибкой интеграции каналов, асинхронное программирование для обеспечения производительности, а также Dependency Injection для связывания компонентов. Применение Docker Compose упрощает развёртывание и обеспечивает воспроизводимость окружения. Тестовое покрытие подтверждает корректность реализации ключевых сценариев использования системы.

На основе проведённой работы можно сделать вывод, что разработанная архитектура соответствует всем требованиям, предъявляемым к современным омниканальным решениям: модульность, масштабируемость, расширяемость и надёжность.

\newpage
\chapter*{Заключение}
\addcontentsline{toc}{chapter}{Заключение}

В ходе прохождения производственной практики я успешно выполнил работы по проектированию и разработке омниканальной платформы для центра обучения, обеспечивающей приём и обработку обращений из различных каналов связи в едином интерфейсе куратора.

Была спроектирована архитектура системы, основанная на принципах Clean Architecture. Определены три уровня: Core (доменные модели и бизнес-логика), API (REST-эндпоинты и WebSocket) и Adapters (реализации PostgreSQL, Telegram, Email, MinIO). Разработанная архитектура обеспечивает слабую связанность компонентов и позволяет в дальнейшем подключать новые каналы связи без изменения существующего кода.

Разработана схема базы данных PostgreSQL, включающая семь таблиц для хранения информации о сообщениях, отправителях, клиентах, каналах клиентов, кураторах и истории назначений. При проектировании использованы внешние ключи, индексы и ограничения для обеспечения целостности данных.

Особое внимание было уделено организации процесса разработки и развёртывания приложения. Для локальной разработки использована контейнеризация Docker и Docker Compose, включающая PostgreSQL 17, MinIO и серверное приложение. Настроены pre-commit хуки, линтеры (ruff) и система статической типизации (mypy). Реализовано тестовое покрытие ключевых сервисов, адаптеров и API-маршрутов с использованием pytest и pytest-asyncio.

В результате прохождения практики были закреплены навыки проектирования архитектуры информационных систем, разработки серверных приложений на Python с использованием асинхронного веб-фреймворка Litestar, проектирования реляционных баз данных PostgreSQL, интеграции внешних API (Telegram Bot API, IMAP/SMTP), разработки клиентских приложений на React и организации развёртывания современных веб-приложений с помощью Docker.

В перспективе система может быть расширена дополнительными каналами связи (SMS, VK), интеграцией с образовательной платформой GetCourse, а также дополнена функциональностью аналитики и отчётности.

\newpage
\appendix
\chapter*{Приложение А \\ Исходный код проекта}
\addcontentsline{toc}{chapter}{Приложение А. Исходный код проекта}

Исходный код разработанного в ходе практики программного комплекса доступен в публичном репозитории на GitHub по адресу:

\medskip
\begin{center}
\url{https://github.com/ProudRykar/PP-2026}
\end{center}
\medskip

Структура репозитория включает следующие основные директории и файлы:

\begin{itemize}
    \item \texttt{app/}~---~серверная часть на Python с использованием фреймворка Litestar
    \begin{itemize}
        \item \texttt{core/}~---~доменная модель, бизнес-логика и порты (интерфейсы)
        \item \texttt{api/}~---~REST-контроллеры, DTO-схемы, обработка исключений
        \item \texttt{adapters/}~---~реализации адаптеров (Telegram, Email, PostgreSQL, MinIO)
    \end{itemize}
    \item \texttt{frontend/}~---~клиентское приложение на React 19 с TypeScript
    \item \texttt{tests/}~---~модульные и интеграционные тесты (pytest, pytest-asyncio)
    \item \texttt{docker-compose.yml}~---~конфигурация развёртывания (PostgreSQL, MinIO, приложение)
    \item \texttt{pyproject.toml}~---~зависимости и настройки инструментов разработки
\end{itemize}

\end{document}
